\documentclass[11pt,a4paper]{article}

\usepackage[utf8]{inputenc}
\usepackage{geometry}
\geometry{margin=2.5cm}
\usepackage{graphicx}
\usepackage{hyperref}
\hypersetup{colorlinks=true, linkcolor=black, urlcolor=blue, citecolor=black, breaklinks=true}
\usepackage{amsmath}
\usepackage{microtype}
\usepackage{seqsplit}
\sloppy

\title{Tagesrhythmus-Vergleich: Übertragung des Carpathian-Bewegungsmodells\\
auf die japanischen Kameradaten}
\author{}
\date{}

\begin{document}
\maketitle
\vspace{-2em}

\section{Ausgangslage}

Die beiden Teilsysteme dieser Arbeit liefern strukturell unterschiedliche
Daten. Die rum\"anischen Carpathian-Daten sind \emph{Trajektorien}: fortlaufende
GPS-Positionen bekannter, einzelner Tiere. Die japanischen Kameradaten sind
dagegen \emph{Detektionen}: Zeitpunkt und Ort einer Ausl\"osung, ohne
Individualkennung und ohne Bewegungspfad zwischen zwei Ausl\"osungen. Eine
Bewegungsanimation im wörtlichen Sinn -- wie sie f\"ur die Carpathian-Daten
umgesetzt wurde -- l\"asst sich aus den japanischen Rohdaten allein nicht
ableiten.

Der Betreuer regte an (Sprechstunde vom 21.08.2026), das in Rum\"anien
gelernte Bewegungsmuster auf die japanischen Kameradaten zu \"ubertragen und
zu pr\"ufen, \emph{"`ob die simulierten B\"aren in \"ahnlichen Zeitr\"aumen wie
echte B\"aren fotografiert werden"'}. Diese Fragestellung wird im Folgenden in
ein konkretes, statistisch auswertbares Verfahren \"ubersetzt.

\section{Datengrundlage und ihre Einschr\"ankungen}

Eine Sichtung der Datenbank (\texttt{bears\_images\_unified}) ergab 52
Bildeintr\"age mit Zeitstempel, verteilt auf 20 Ger\"ate -- nach Ausschluss
zweier Platzhalter-Ger\"ate-IDs (\texttt{199999}, \texttt{999999}), die
vermutlich aus einer fehlgeschlagenen Ger\"ate-Erkennung beim
E-Mail-Import stammen. Keine einzelne Kamera erreicht mehr als elf
Zeitstempel. Eine Auswertung \emph{pro Kamera} ist damit nicht sinnvoll
durchf\"uhrbar; die Analyse muss \"uber alle Kameras gepoolt erfolgen, was die
r\"aumliche Aufl\"osung der Aussage reduziert, aber weiterhin einen zeitlichen
Vergleich erlaubt.

Ein erster Blick auf die gepoolte Verteilung zeigt ein auff\"alliges Muster:
90\,\% aller Ausl\"osungen liegen zwischen 18 und 6 Uhr. Asiatische
Kragenb\"aren (\emph{Ursus thibetanus}) gelten in der Literatur eigentlich als
cathemeral, also grunds\"atzlich sowohl tag- als auch nachtaktiv. Die nahezu
ausschlie{\ss}liche Nachtaktivit\"at in den vorliegenden Daten ist damit selbst
ein bemerkenswerter, eigenst\"andig zeigbarer Befund -- konsistent mit der
in mehreren Studien beschriebenen Verschiebung der Aktivit\"at in die Nacht
unter erh\"ohtem menschlichem Nutzungsdruck, was angesichts der Lage der
Ger\"ate an Abschreckungsanlagen (also an Orten mit Mensch-Wildtier-Konflikt)
plausibel ist.

\section{Methodik}

\subsection{\"Uberlappungskoeffizient}

Zum Vergleich zweier Tagesrhythmus-Verteilungen wird der
\"Uberlappungskoeffizient $\Delta$ verwendet, das Standardma{\ss} der
Kamerafallen-Forschung f\"ur diesen Zweck:
\begin{equation}
  \Delta = \int_0^{24} \min\bigl(f_1(t),\, f_2(t)\bigr)\, dt
\end{equation}
mit $f_1, f_2$ als (auf $\int f = 1$ normierten) Dichtefunktionen der
Ausl\"osezeiten zweier Gruppen. $\Delta = 0$ bedeutet keinerlei zeitliche
\"Uberschneidung, $\Delta = 1$ identische Muster. Das Ma{\ss} und die
zugeh\"orige Sch\"atzmethodik wurden von Ridout und Linkie eingef\"uhrt und
sind im weit verbreiteten R-Paket \texttt{overlap} implementiert.\footnote{Ridout,
M.S., Linkie, M.: \emph{Estimating overlap of daily activity patterns from
camera trap data}. Journal of Agricultural, Biological and Environmental
Statistics (2009). Sowie: Meredith, M., Ridout, M.: \emph{Overview of the
overlap package}, CRAN-Vignette.}

In dieser Arbeit wird $\Delta$ nicht \"uber die von-Mises-Kern-Implementierung
des R-Pakets, sondern \"uber eine selbst implementierte, gespiegelte
("`wrapped"') Gau{\ss}-Kerndichtesch\"atzung auf dem 24-Stunden-Kreis
berechnet (Datei \texttt{activity\_overlap.py}, Funktion
\texttt{circular\_kde}). Die Spiegelung stellt sicher, dass z.\,B. 23:30 Uhr
und 00:30 Uhr korrekt als benachbart behandelt werden, nicht als
entgegengesetzte Enden einer Zeitachse. Die Wahl dieser vereinfachten,
selbst kontrollierten Implementierung wurde getroffen, um keine zus\"atzliche
R-Abh\"angigkeit in ein sonst rein Python-basiertes System einzuf\"uhren; das
zugrunde liegende Prinzip ist identisch.

\subsection{Simulierte Vergleichsgruppe: Monte-Carlo-\"Ubertragung}

Als Vergleichsgruppe zu den echten Detektionen dient keine einzelne
Simulation, sondern eine \emph{Monte-Carlo-Stichprobe} aus 300 unabh\"angigen
simulierten B\"arenpfaden (je 48\,h), jeweils gestartet an einem
zuf\"allig gew\"ahlten echten japanischen Kamerastandort. Geschlecht, Alter
und Saison werden pro Lauf zuf\"allig gezogen, um \"uber die unbekannte
Demografie der japanischen Population zu mitteln. Die Bewegungslogik selbst
ist identisch mit der bereits f\"ur den Realit\"atscheck der Simulation
implementierten Schrittgenerierung (empirische, stundenaufgel\"oste
Schrittweiten aus den rum\"anischen \texttt{BearSpeeds}-Daten, zuf\"allige
Richtung). N\"ahert sich ein simulierter Pfadpunkt einer beliebigen echten
Kamera auf $\leq 150\,$m, wird dies als \emph{virtuelle Ausl\"osung}
gez\"ahlt und deren Uhrzeit gespeichert.

Dieses Vorgehen -- Bewegung simulieren, virtuelle Kamerafallen auswerten,
mit echten Kamerafallendaten vergleichen -- ist ein etabliertes Studiendesign
in der Wildtier\"okologie und wurde u.\,a. zur Validierung von
Kamerafallen-Dichtesch\"atzmethoden anhand simulierter Tierpopulationen
eingesetzt.\footnote{Vgl. z.\,B. Calenge, C. et al. sowie Santini, L. et al.
zu simulationsbasierten Kamerafallen-Studiendesigns; genaue Quellenangaben
sind \"uber die Recherche-Historie des Bachelorarbeitsprojekts nachvollziehbar
und vor \"Ubernahme in das Literaturverzeichnis im Volltext zu verifizieren.}

\subsection{Unsicherheitsangabe}

Da die reale Stichprobe klein ist (52 Detektionen, nach Bereinigung), wird
$\Delta$ nicht als einzelner, unbegr\"undet pr\"azise wirkender Wert
berichtet, sondern mit einem 95\,\%-Bootstrap-Konfidenzintervall (1000
Wiederholungen, Ziehen mit Zur\"ucklegen aus beiden Stichproben). Ein
Beispiel-Ergebnis (mit einer synthetischen, dem realen Muster nachempfundenen
Stichprobe erzeugt, siehe Abbildung~\ref{fig:overlap}) liefert $\Delta =
0{,}59$ bei einem Konfidenzintervall von $[0{,}51,\ 0{,}64]$.

\begin{figure}[h]
  \centering
  \includegraphics[width=0.85\textwidth]{activity_overlap_beispiel.png}
  \caption{Beispielhafter Tagesrhythmus-Vergleich: echte Detektionsdichte
    (blau) gegen simulierte Detektionsdichte (orange), graue Fl\"ache zeigt
    die \"Uberlappung. Erzeugt mit einer synthetischen Stichprobe zur
    Illustration der Methodik -- vor Verwendung in der Arbeit mit dem
    tats\"achlichen Datenbank-Ergebnis zu ersetzen.}
  \label{fig:overlap}
\end{figure}

\section{Beitrag zur Arbeit}

Diese Analyse leistet unabh\"angig vom konkreten $\Delta$-Wert mehrere
Beitr\"age:

\begin{enumerate}
  \item Sie beantwortet die im Meeting gestellte Frage (zeitliche
        \"Ubereinstimmung simulierter und echter B\"arenbewegung) mit einer
        einzelnen, interpretierbaren Kennzahl statt einer subjektiven
        Einsch\"atzung.
  \item Sie macht die Grenzen der Datenlage explizit sichtbar (kleine
        Stichprobe, keine Kamera-individuelle Auswertung m\"oglich) statt sie
        zu verschweigen -- und liefert trotzdem ein verwertbares, statistisch
        abgesichertes Ergebnis.
  \item Ein niedriger $\Delta$-Wert w\"are \emph{kein} Fehlschlag der
        Umsetzung, sondern eine valide wissenschaftliche Aussage: dass
        rum\"anische Braunb\"ar-Bewegungsmuster die Aktivit\"at japanischer
        Kragenb\"aren zeitlich nicht erkl\"aren -- plausibel angesichts der
        unterschiedlichen Art und \"Okologie.
  \item Der Ansatz ist erweiterbar: Ein Vergleichslauf mit publizierten
        japanischen Bewegungsparametern (statt der rum\"anischen) anstelle der
        \texttt{BearSpeeds}-Tabelle w\"urde direkt zeigen, wie stark der
        Artunterschied den $\Delta$-Wert beeinflusst -- ein m\"oglicher
        n\"achster Ausbauschritt.
\end{enumerate}

\section{Grenzen}

Die Aussagekraft ist durch die geringe Stichprobengr\"o{\ss}e (52 Detektionen)
eingeschr\"ankt; das Ergebnis ist explorativ zu lesen, nicht als
konfirmatorischer Test. Die gepoolte Auswertung verdeckt m\"ogliche
Unterschiede zwischen einzelnen Kamerastandorten (z.\,B. unterschiedliche
Habitatn\"ahe). Die Simulation nutzt weiterhin ein an rum\"anischen
Braunb\"aren kalibriertes Bewegungsmodell fu\"r eine andere Art
(\emph{Ursus thibetanus} statt \emph{Ursus arctos}) -- eine bewusste, im
Kontext der Aufgabenstellung nachvollziehbare Vereinfachung, die als
Limitation zu benennen ist.

\end{document}
